\documentclass[11pt]{article}
\usepackage{amsmath,amssymb,amsthm}
\usepackage[margin=1.1in]{geometry}
\usepackage{booktabs}

\newtheorem{proposition}{Proposition}
\newtheorem{remark}{Remark}

\title{Activation Function Derivatives:\\
Every Nonlinearity in the Repository, Differentiated From Scratch}
\author{Yuval Lavie}
\date{\today}

\begin{document}
\maketitle

\noindent
Backpropagation multiplies each layer's upstream gradient by the local
derivative of its activation. This page derives every one of those
local derivatives explicitly --- no step skipped --- citing the rules
page (\texttt{rules.tex}) for the calculus. A recurring theme:
the good activations have derivatives expressible \emph{in terms of
their own outputs}, which is what makes the backward pass cheap ---
the forward value is already cached.

\section{Sigmoid}

\begin{equation}
  \sigma(x) = \frac{1}{1 + e^{-x}}.
  \label{eq:sigmoid}
\end{equation}

\paragraph{Derivation.} Write $\sigma(x) = (1 + e^{-x})^{-1}$ and apply
the chain rule twice (power rule on the outer function, then
$\frac{d}{dx}e^{-x} = -e^{-x}$ on the inner):
\begin{align*}
  \sigma'(x)
  &= -\,(1 + e^{-x})^{-2} \cdot \frac{d}{dx}\bigl(1 + e^{-x}\bigr)
   = -\,(1 + e^{-x})^{-2} \cdot \bigl(-e^{-x}\bigr)
   = \frac{e^{-x}}{(1 + e^{-x})^{2}}.
\end{align*}
Now factor into a product of two sigmoids:
\begin{align*}
  \frac{e^{-x}}{(1+e^{-x})^{2}}
  &= \frac{1}{1+e^{-x}} \cdot \frac{e^{-x}}{1+e^{-x}}
   = \sigma(x)\cdot\frac{(1 + e^{-x}) - 1}{1+e^{-x}}
   = \sigma(x)\bigl(1 - \sigma(x)\bigr).
\end{align*}
\begin{equation}
  \boxed{\;\sigma'(x) = \sigma(x)\,\bigl(1 - \sigma(x)\bigr)\;}
  \label{eq:sigmoidprime}
\end{equation}

\begin{remark}[Reading \eqref{eq:sigmoidprime}]
The derivative is the output times its complement --- a Bernoulli
variance $p(1-p)$. It peaks at $\sigma'(0) = 1/4$ and decays
exponentially in $|x|$: this $\le 1/4$ ceiling, compounded across
layers, is the classical \emph{vanishing gradient} problem, and the
reason sigmoid retired from hidden layers in favor of ReLU.
Two more useful identities: $\sigma(-x) = 1 - \sigma(x)$, and by the
inverse function rule the logit
$\operatorname{logit}(p) = \ln\frac{p}{1-p}$ has derivative
$\frac{1}{p(1-p)}$.
\end{remark}

\section{Hyperbolic tangent}

\begin{equation}
  \tanh(x) = \frac{e^{x} - e^{-x}}{e^{x} + e^{-x}}
           = \frac{\sinh x}{\cosh x}.
  \label{eq:tanh}
\end{equation}

\paragraph{Derivation.} Quotient rule with $f = \sinh$, $g = \cosh$,
using $\sinh' = \cosh$, $\cosh' = \sinh$, and the identity
$\cosh^2 - \sinh^2 = 1$:
\begin{align*}
  \tanh'(x)
  = \frac{\cosh x \cdot \cosh x - \sinh x \cdot \sinh x}{\cosh^{2}x}
  = \frac{\cosh^2 x - \sinh^2 x}{\cosh^{2}x}
  = \frac{1}{\cosh^{2}x} = 1 - \frac{\sinh^2 x}{\cosh^2 x}.
\end{align*}
\begin{equation}
  \boxed{\;\tanh'(x) = 1 - \tanh^{2}(x)\;}
  \label{eq:tanhprime}
\end{equation}

\begin{remark}
Again the derivative is a function of the output alone. Since
$\tanh(x) = 2\sigma(2x) - 1$, tanh is a rescaled sigmoid; its
derivative peaks at $\tanh'(0) = 1$, four times sigmoid's peak, and its
outputs are zero-centered --- both reasons it trains better than
sigmoid in classical (pre-ReLU) networks.
\end{remark}

\section{ReLU and its leaky variants}

\begin{equation}
  \mathrm{ReLU}(x) = \max(0, x)
  = \begin{cases} x & x > 0\\ 0 & x \le 0. \end{cases}
  \label{eq:relu}
\end{equation}

\paragraph{Derivation.} On each open half-line the function is affine,
so differentiate piecewise:
\begin{equation}
  \mathrm{ReLU}'(x) =
  \begin{cases}
    1 & x > 0\\
    0 & x < 0,
  \end{cases}
  \qquad
  \partial\,\mathrm{ReLU}(0) = [0, 1].
  \label{eq:reluprime}
\end{equation}
At $x = 0$ the two one-sided slopes ($0$ and $1$) disagree, so no
derivative exists; ReLU is convex, so the subdifferential $[0,1]$
(rules page) stands in, and frameworks pick the element $0$. The
backward pass is a mask: gradients flow unchanged through active units
and are killed at inactive ones --- no shrinking factor, hence no
vanishing gradient on the active path, but also \emph{dead units} if a
unit's input stays negative forever.

\paragraph{Leaky ReLU / PReLU.} Replace the zero slope with
$\alpha > 0$ (fixed for Leaky, learned for PReLU):
\begin{equation}
  f(x) = \begin{cases} x & x > 0\\ \alpha x & x \le 0 \end{cases}
  \qquad\Longrightarrow\qquad
  f'(x) = \begin{cases} 1 & x > 0\\ \alpha & x < 0, \end{cases}
  \label{eq:leaky}
\end{equation}
and for PReLU the gradient with respect to the learned slope is
$\partial f / \partial \alpha = x\,\mathbf{1}[x \le 0]$.

\section{ELU}

For $\alpha > 0$:
\begin{equation}
  \mathrm{ELU}(x) =
  \begin{cases}
    x & x > 0\\
    \alpha\,(e^{x} - 1) & x \le 0.
  \end{cases}
  \label{eq:elu}
\end{equation}

\paragraph{Derivation.} The positive branch is the identity; on the
negative branch, linearity and $\frac{d}{dx}e^x = e^x$ give
$\alpha e^{x}$, which equals $\mathrm{ELU}(x) + \alpha$:
\begin{equation}
  \mathrm{ELU}'(x) =
  \begin{cases}
    1 & x > 0\\
    \alpha e^{x} = \mathrm{ELU}(x) + \alpha & x \le 0.
  \end{cases}
  \label{eq:eluprime}
\end{equation}
At $x=0$ both branches give value $0$ and slope $1$
($\alpha e^0 = \alpha$ only matches when $\alpha = 1$; for
$\alpha \ne 1$ the derivative jumps from $\alpha$ to $1$). Unlike
ReLU, saturating units still receive gradient $\alpha e^x > 0$.

\section{Softplus}

\begin{equation}
  \zeta(x) = \ln\bigl(1 + e^{x}\bigr).
  \label{eq:softplus}
\end{equation}

\paragraph{Derivation.} Log-derivative rule with $f(x) = 1 + e^{x}$:
\begin{equation}
  \zeta'(x) = \frac{e^{x}}{1 + e^{x}}
  = \frac{1}{e^{-x} + 1}
  \;\Longrightarrow\;
  \boxed{\;\zeta'(x) = \sigma(x)\;}
  \label{eq:softplusprime}
\end{equation}
Softplus is the antiderivative of the sigmoid --- the smooth ReLU
($\zeta(x) \to \max(0,x)$ pointwise as its temperature sharpens). This
pairing (softplus, sigmoid) is the same (loss, gradient) pairing that
makes the logistic loss's gradient a sigmoid on the losses page.

\section{SiLU / Swish}

\begin{equation}
  \mathrm{silu}(x) = x\,\sigma(x).
  \label{eq:silu}
\end{equation}

\paragraph{Derivation.} Product rule with
$\sigma' = \sigma(1-\sigma)$ from \eqref{eq:sigmoidprime}:
\begin{align*}
  \mathrm{silu}'(x)
  = \sigma(x) + x\,\sigma(x)\bigl(1 - \sigma(x)\bigr)
  = \sigma(x)\bigl(1 + x - x\,\sigma(x)\bigr).
\end{align*}
\begin{equation}
  \boxed{\;\mathrm{silu}'(x)
  = \sigma(x) + x\,\sigma(x)\bigl(1-\sigma(x)\bigr)
  = \mathrm{silu}(x) + \sigma(x)\bigl(1 - \mathrm{silu}(x)\bigr)\;}
  \label{eq:siluprime}
\end{equation}
(The second form is output-only, again.) SiLU is smooth, non-monotone
(a small dip for $x \approx -1.28$), and its derivative can be
slightly negative --- gradients can flow ``backwards'' through mildly
negative inputs, unlike ReLU's hard zero.

\section{GELU}

With $\Phi$ the standard normal CDF and
$\varphi(x) = \frac{1}{\sqrt{2\pi}}e^{-x^2/2}$ its density:
\begin{equation}
  \mathrm{GELU}(x) = x\,\Phi(x).
  \label{eq:gelu}
\end{equation}

\paragraph{Derivation.} Product rule, using the fundamental theorem of
calculus for $\Phi' = \varphi$:
\begin{equation}
  \boxed{\;\mathrm{GELU}'(x) = \Phi(x) + x\,\varphi(x)\;}
  \label{eq:geluprime}
\end{equation}
GELU is ``expected ReLU under Gaussian gating'':
$x \cdot \Pr(Z \le x)$ for $Z \sim \mathcal{N}(0,1)$. The common
tanh approximation
$\mathrm{GELU}(x) \approx \tfrac{x}{2}\bigl(1 +
\tanh\bigl[\sqrt{2/\pi}\,(x + 0.044715\,x^{3})\bigr]\bigr)$
differentiates mechanically by the product and chain rules with
\eqref{eq:tanhprime}; the transformer notes (\texttt{llm/}) use the
exact form.

\section{Softmax: the Jacobian}

Softmax maps logits $z \in \mathbb{R}^{K}$ to a distribution:
\begin{equation}
  s_i(z) = \frac{e^{z_i}}{\sum_{k=1}^{K} e^{z_k}}.
  \label{eq:softmax}
\end{equation}
It is vector-valued, so its derivative is a $K \times K$ Jacobian
$\partial s_i / \partial z_j$, with two cases.

\paragraph{Case $i = j$.} Quotient rule with $f = e^{z_i}$,
$g = \Sigma := \sum_k e^{z_k}$, noting
$\partial \Sigma / \partial z_i = e^{z_i}$:
\begin{align*}
  \frac{\partial s_i}{\partial z_i}
  = \frac{e^{z_i}\,\Sigma - e^{z_i}\,e^{z_i}}{\Sigma^{2}}
  = \frac{e^{z_i}}{\Sigma}\cdot\frac{\Sigma - e^{z_i}}{\Sigma}
  = s_i\,(1 - s_i).
\end{align*}

\paragraph{Case $i \ne j$.} Now the numerator $e^{z_i}$ is a constant
with respect to $z_j$, so only the denominator moves; by the chain rule
on $\Sigma^{-1}$:
\begin{align*}
  \frac{\partial s_i}{\partial z_j}
  = e^{z_i} \cdot \Bigl(-\Sigma^{-2}\Bigr)\cdot e^{z_j}
  = -\,\frac{e^{z_i}}{\Sigma}\cdot\frac{e^{z_j}}{\Sigma}
  = -\,s_i\,s_j.
\end{align*}

\paragraph{Both cases at once.} With the Kronecker delta
$\delta_{ij}$:
\begin{equation}
  \boxed{\;\frac{\partial s_i}{\partial z_j}
  = s_i\,\bigl(\delta_{ij} - s_j\bigr)\;}
  \qquad\text{i.e.}\qquad
  J_s = \operatorname{diag}(s) - s\,s^{\top}.
  \label{eq:softmaxjac}
\end{equation}

\begin{remark}[Reading \eqref{eq:softmaxjac}]
The diagonal entries $s_i(1-s_i)$ are exactly sigmoid's
\eqref{eq:sigmoidprime} --- softmax with $K=2$ \emph{is} a sigmoid of
the logit difference. The off-diagonal $-s_i s_j$ encodes competition:
raising one logit necessarily lowers every other probability. The
Jacobian is symmetric, positive semidefinite, and singular ---
$J_s \mathbf{1} = 0$, because shifting all logits by a constant leaves
softmax unchanged (the invariance the log-sum-exp trick exploits for
stability). In practice one never materializes $J_s$: the losses page
shows that composing it with cross-entropy collapses the whole product
to $s - y$.
\end{remark}

\section{Summary table}

Derivatives in terms of the \emph{output} $a = f(x)$ where possible:

\begin{center}
\begin{tabular}{llll}
\toprule
Activation & $f(x)$ & $f'$ & Range of $f'$\\
\midrule
Sigmoid & $1/(1+e^{-x})$ & $a(1-a)$ & $(0, \tfrac14]$\\
Tanh & $\tanh x$ & $1 - a^2$ & $(0, 1]$\\
ReLU & $\max(0,x)$ & $\mathbf{1}[x>0]$ & $\{0, 1\}$\\
Leaky ReLU & $\max(\alpha x, x)$ & $\alpha$ or $1$ & $\{\alpha, 1\}$\\
ELU & eq.\ \eqref{eq:elu} & $1$ or $a + \alpha$ & $(0, 1]$\\
Softplus & $\ln(1+e^x)$ & $\sigma(x)$ & $(0,1)$\\
SiLU & $x\,\sigma(x)$ & $a + \sigma(x)(1-a)$ & $\approx(-0.09, 1.1)$\\
GELU & $x\,\Phi(x)$ & $\Phi(x) + x\varphi(x)$ & $\approx(-0.13, 1.13)$\\
Softmax & eq.\ \eqref{eq:softmax} & $s_i(\delta_{ij} - s_j)$ &
Jacobian\\
\bottomrule
\end{tabular}
\end{center}

\medskip
\noindent
The pattern to remember: sigmoid, tanh, ELU, and SiLU all have
output-expressible derivatives (cache the forward value, get the
backward almost free); ReLU's derivative is a mask; softmax's is a
rank-one correction of a diagonal. The loss-functions page composes
these with the loss derivatives to produce the clean gradients
($\sigma(z) - y$, $s - y$) that make logistic and softmax regression
train so gracefully.

\end{document}
